\documentclass[a4paper]{article}
\usepackage[no-math,cm-default]{fontspec}
\usepackage[amsbb,subscriptcorrection,zswash,mtpcal,mtphrb]{mtpro2}
\usepackage{xunicode}
\usepackage{xgreek}
\usepackage{amsmath}
\RequirePackage[left=3.00cm, right=3.00cm, top=2.00cm, bottom=3.00cm]{geometry}
\defaultfontfeatures{Mapping=tex-text,Scale=MatchLowercase}
\setmainfont[Mapping=tex-text,Numbers=Lining,Scale=1.0,BoldFont={Times New Roman Bold}]{Times New Roman}

%------TIKZ - ΣΧΗΜΑΤΑ - ΓΡΑΦΙΚΕΣ ΠΑΡΑΣΤΑΣΕΙΣ ----
\usepackage{tikz,pgfplots}
\usepackage{tkz-euclide}
\usetkzobj{all}

%------------- ΔΙΑΦΟΡΑ ----------
\usepackage{etoolbox}
\usepackage{calc}
\usepackage{hhline}
\usepackage{graphicx}
\usepackage{multicol}
\usepackage{multirow}
\usepackage{enumitem}
\usepackage{tabularx}
\usepackage{systeme}
%-------- ΜΑΘΗΜΑΤΙΚΑ ΕΡΓΑΛΕΙΑ ---------
\usepackage{mathtools}
%----------------------
%-------- ΠΙΝΑΚΕΣ ---------
\usepackage{booktabs}
%----------------------
%----- ΥΠΟΛΟΓΙΣΤΗΣ ----------
\usepackage{calculator}
%----------------------------

%---------- ΛΙΣΤΕΣ ----------------------
\newlist{alist}{enumerate}{3}
\setlist[alist]{itemsep=0mm,label=\alph*.}
\newlist{balist}{enumerate}{3}
\setlist[balist]{itemsep=0mm,label=\bf\alph*.}
\newlist{Alist}{enumerate}{3}
\setlist[Alist]{itemsep=0mm,label=\Alph*.}
\newlist{bAlist}{enumerate}{3}
\setlist[bAlist]{itemsep=0mm,label=\bf\Alph*.}
\newlist{rlist}{enumerate}{3}
\setlist[rlist]{itemsep=0mm,label=\roman*.}
\newlist{brlist}{enumerate}{3}
\setlist[brlist]{itemsep=0mm,label=\bf\roman*.}

%------- Αρίθμιση με ελληνικά γράμματα ----------
\makeatletter
\renewrobustcmd{\anw@true}{\let\ifanw@\iffalse}
\renewrobustcmd{\anw@false}{\let\ifanw@\iffalse}\anw@false
\newrobustcmd{\noanw@true}{\let\ifnoanw@\iffalse}
\newrobustcmd{\noanw@false}{\let\ifnoanw@\iffalse}\noanw@false
\renewrobustcmd{\anw@print}{\ifanw@\ifnoanw@\else\numer@lsign\fi\fi}
\makeatother
\usepackage{ntheorem}
\newtheorem{theorem}{Θεώρημα}
\usepackage{amsmath}
\DeclareMathOperator{\Res}{Res}
\theoremstyle{definition}
\newtheorem{definition}{Ορισμός}


\begin{document}

\begin{definition}[Σειρά Fourier]
Έστω $ E $ ο χώρος των κατά τμήματα συνεχών συναρτήσεων με πεδίο ορισμού το $ [-\pi,\pi] $ και $ f\in E $.H ακόλουθη σειρά ονομάζεται \textbf{σειρά Fourier} της $ f $ και δίνεται από τη σχέση
\[ f(x)=\frac{a_0}{2}+\sum_{n=1}^{\infty}{(a_n\cos{(nx)}+b_n\sin{(nx)})} \]
όπου $ a_n,b_n $ με $ n=1,2,\ldots $ είναι οι συντελεστές Fourier και δίνονται από τους τύπους:
\[ a_n=\frac{1}{\pi}\int_{-\pi}^{\pi}{f(x)\cos{(nx)}}d x\ \ \textrm{και}\ \ b_n=\frac{1}{\pi}\int_{-\pi}^{\pi}{f(x)\sin{(nx)}}d x \]
\end{definition}

\begin{theorem}[Θεώρημα ολοκληρωτικών υπολοίπων]
Έστω $ f $ είναι ολόμορφη πάνω και στο εσωτερικό μιας απλής κλειστής θετικά προσανατολισμένης τμηματικά λείας καμπύλης $ \gamma $ με εξαίρεση ένα πεπερασμένο πλήθος απομονωμένων ανωμάλων σημείων $ z_1,z_2\ldots,z_n $ στο εσωτερικό της $ \gamma $. Τότε
\[
  \frac{1}{2\pi i}\oint_\gamma\! f(z)dz = \sum_{k=1}^n\Res(f;z_k)\,.
\]
\end{theorem}

\begin{theorem}[Θεώρημα απόκλισης του Gauss]
Έστω $ \Omega $ ένα υποσύνολο του $ \mathbb{R}^3 $ κλειστό και φραγμένο με κατά τμήματα $ C^2 $ σύνορο $ \partial\Omega $. Τότε αν η $ f $ με πεδίο τιμών το $ \mathbb{R}^3 $ είναι $ C^1 $ συνάρτηση στο $ \Omega $ τότε ισχύει (αλλαγή)
\[
 \iiint_{\Omega}{\nabla f(x,y,z)dxdydz}=\oiint_{\partial\Omega}{f(x,y,z)\cdot n(x,y,z)d\sigma}
\]
\end{theorem}
\end{document}
